\section{A uniform Laguerre estimate}\label{app:laguerre}

For $\alpha=1/2$, use the normalized Laguerre function
\[
\mathcal L_m^{(1/2)}(r)
 =a_mL_m^{(1/2)}(r)e^{-r/2}r^{1/4},
\qquad
a_m=\sqrt{\frac{m!}{\Gamma(m+3/2)}},
\qquad
\nu=4m+3.
\]
Proposition~3.2 of Imekraz, Robert, and Thomann~\cite{ImekrazRobertThomann2016} gives constants $C,c>0$, independent of $m$ and $r$, such that
\[
|\mathcal L_m^{(1/2)}(r)|\le C(r\nu)^{1/4}
\quad (0\le r\le\nu^{-1}),
\]
\[
|\mathcal L_m^{(1/2)}(r)|\le C(r\nu)^{-1/4}
\quad (\nu^{-1}\le r\le\nu/2),
\]
\[
|\mathcal L_m^{(1/2)}(r)|
 \le C\nu^{-1/4}\bigl(\nu^{1/3}+|\nu-r|\bigr)^{-1/4}
 \quad (\nu/2\le r\le3\nu/2),
\]
and
\[
|\mathcal L_m^{(1/2)}(r)|\le Ce^{-cr}
\quad (r\ge3\nu/2).
\]
Since
\[
Q_m(r):=\sqrt r\,e^{-r/2}|L_m^{(1/2)}(r)|
 =a_m^{-1}r^{1/4}|\mathcal L_m^{(1/2)}(r)|,
\qquad
a_m^{-1}\asymp m^{1/4}\asymp\nu^{1/4},
\]
the four regions can be checked without hiding the powers of $m$.

In Region I, $0\le r\le\nu^{-1}$, the estimate gives
\[
Q_m(r)
 \le C\nu^{1/4}r^{1/4}(r\nu)^{1/4}
 =C(r\nu)^{1/2}\le C.
\]
In Region II, $\nu^{-1}\le r\le\nu/2$, it gives
\[
Q_m(r)
 \le C\nu^{1/4}r^{1/4}(r\nu)^{-1/4}
 =C.
\]
In the turning region, $\nu/2\le r\le3\nu/2$, we obtain
\[
Q_m(r)
 \le C r^{1/4}\bigl(\nu^{1/3}+|\nu-r|\bigr)^{-1/4}
 \le C\nu^{1/6}.
\]
Multiplication by the remaining factor in the desired envelope gives
\[
e^{-r/6}Q_m(r)
 \le C\nu^{1/6}e^{-\nu/12}\le C.
\]
Finally, in the outer region $r\ge3\nu/2$, the exponential estimate yields
\[
Q_m(r)
 \le C\nu^{1/4}r^{1/4}e^{-cr}
 \le C\,\operatorname{poly}(r,\nu)e^{-cr}.
\]
After multiplication by $e^{-r/6}$ this is uniformly bounded. Hence
\[
\boxed{
\sup_{m\ge1,\,r\ge0}\sqrt r\,e^{-2r/3}|L_m^{(1/2)}(r)|<\infty.
}
\]
The case $m=0$ is immediate from $L_0^{(1/2)}(r)=1$. Substituting $r=3x^2/2$ converts this bound into the constant $C_L$ used in Section~\ref{sec:sharpness}.
